\bookchapter{结语：保持好奇，保持清醒}

恭喜你读到了这里。

这本书从手机硬件拆到软件生态，从日常设置挖到开发者模式，从安全防护聊到数字遗产。如果它帮你解决了一个实际问题——哪怕只是学会了关掉一个不该开的权限、养成了一次备份的习惯——那这些字就没白写。

但比任何具体技巧都重要的，是两件事：

\begin{itemize}
  \item \textbf{保持好奇}：手机每年都在变，AI在变，折叠屏在变，卫星通信在变。今天学的操作步骤明天可能就过时了，但"想知道为什么"的好奇心永远不会过时。多问一句"这个权限是干什么的""这个数据存在了哪里""这个功能会不会泄露隐私"，比记住一百个操作步骤都有用。
  \item \textbf{保持清醒}：手机是工具，不是主人。它设计出来就是为了吸引你的注意力、占据你的时间、引导你的消费。用它是你的选择，被它用是你的失控。定期抬头看看屏幕时间，定期想想"我是在用手机做事，还是在被手机消磨"，这不是矫情，这是数字时代的生存本能。
\end{itemize}

技术在变，但人的需求没变——我们需要安全感、需要连接、需要被尊重。手机只是一个载体。真正的“手机课”不在书里，在你每天解锁屏幕的那一瞬间。

\begin{tcbpractice}{最后一个练习}
合上这本书，拿起你的手机，做一件事：\textbf{用30秒想想，你的手机里有哪三样东西是“丢了会真正心疼”的？}然后，确保它们被备份了。这本书就值了。

\end{tcbpractice}
